\documentclass[11pt,letterpaper]{article}
\usepackage[margin=1in]{geometry}
\usepackage{booktabs}
\usepackage{longtable}
\usepackage{array}
\usepackage{enumitem}
\usepackage{xcolor}
\usepackage{titlesec}
\usepackage{fancyhdr}
\usepackage{hyperref}
\usepackage{colortbl}
\usepackage{tabularx}
\usepackage{makecell}

% Colors
\definecolor{headerblue}{RGB}{47,68,117}
\definecolor{week1color}{RGB}{139,90,43}
\definecolor{week2color}{RGB}{47,100,71}
\definecolor{week3color}{RGB}{130,50,50}
\definecolor{lightgray}{RGB}{245,245,245}
\definecolor{accentblue}{RGB}{70,100,160}

% Header/Footer
\pagestyle{fancy}
\fancyhf{}
\fancyhead[L]{\small\textcolor{headerblue}{MUS 262 --- Three Eras Compared}}
\fancyhead[R]{\small\textcolor{headerblue}{Quiz Study Guide}}
\fancyfoot[C]{\thepage}
\renewcommand{\headrulewidth}{0.5pt}

% Section formatting
\titleformat{\section}{\Large\bfseries\color{headerblue}}{}{0em}{}[\vspace{-0.5em}\rule{\textwidth}{0.5pt}]
\titleformat{\subsection}{\large\bfseries\color{accentblue}}{}{0em}{}
\titleformat{\subsubsection}{\normalsize\bfseries}{}{0em}{}

\hypersetup{colorlinks=true, linkcolor=headerblue, urlcolor=accentblue}

\setlength{\parindent}{0pt}
\setlength{\parskip}{0.5em}

\begin{document}

\begin{center}
{\Huge\bfseries\textcolor{headerblue}{Three Eras Compared}}\\[0.3em]
{\Large Quiz Study Guide}\\[0.5em]
{\large MUS/AAS 262 --- Jazz History: Many Sounds, Many Voices}\\[0.3em]
{\normalsize Early Jazz $\cdot$ Big Bands \& Swing $\cdot$ Bebop}
\end{center}

\vspace{1em}

\section{The One-Sentence Distinction}

\begin{itemize}[leftmargin=1.5em]
  \item \textbf{\textcolor{week1color}{Week 1: Early Jazz (1900--1928)}} --- Jazz is \emph{born}: small groups improvise \emph{together} in New Orleans.
  \item \textbf{\textcolor{week2color}{Week 2: Big Bands \& Swing (1925--1945)}} --- Jazz goes \emph{mainstream}: large orchestras make America \emph{dance}.
  \item \textbf{\textcolor{week3color}{Week 3: Bebop (1940--1950)}} --- Jazz becomes \emph{art}: small combos play fast, complex music you \emph{listen} to, not dance to.
\end{itemize}

\section{Side-by-Side Comparison}

\renewcommand{\arraystretch}{1.4}
\begin{small}
\begin{longtable}{>{\bfseries}p{2.6cm} p{3.8cm} p{3.8cm} p{3.8cm}}
\toprule
& \textbf{\textcolor{week1color}{Early Jazz}} & \textbf{\textcolor{week2color}{Swing / Big Bands}} & \textbf{\textcolor{week3color}{Bebop}} \\
\midrule
\endfirsthead
\toprule
& \textbf{\textcolor{week1color}{Early Jazz}} & \textbf{\textcolor{week2color}{Swing / Big Bands}} & \textbf{\textcolor{week3color}{Bebop}} \\
\midrule
\endhead
Dates & 1900--1928 & 1925--1945 & 1940--1950 \\
\midrule
Birthplace & New Orleans & Harlem (Cotton Club), Kansas City & Harlem (Minton's Playhouse, Clark Monroe's) \\
\midrule
Ensemble size & Small group (5--7) & Big band (15 pieces) & Small combo (4--6) \\
\midrule
Core instruments & Trumpet/cornet, clarinet, trombone & 3 trumpets, 3 trombones, 5 saxes, 4 rhythm & Trumpet/sax, piano, bass, drums \\
\midrule
Dominant wind & Clarinet \& trumpet & Saxophone replaces clarinet & Alto sax (Parker), trumpet (Dizzy) \\
\midrule
Improvisation & Collective (heterophony) --- everyone at once & Solos within arranged frameworks & Extended virtuosic solos over complex chord changes \\
\midrule
Tempo & Moderate & Moderate, danceable & Very fast (or very slow ballads) \\
\midrule
Rhythm feel & Swing rhythm emerging; beats 2 \& 4 & Four-on-the-floor (bass drum/hi-hat); smooth swing & Drummer on ride cymbal; pushing forward on beat \\
\midrule
Purpose & Community music, funerals, dances & Dance music --- jazz IS pop music & Art music / listening --- deliberately NOT for dancing \\
\midrule
Audience & Local New Orleans & Mass national audience (radio, film) & Small insider audience; counterculture \\
\midrule
Harmonic complexity & Blues scale, simple harmonies & 32-bar AABA standards, rhythm changes & Upper extensions (9ths, 11ths, 13ths); ``playing the changes'' \\
\midrule
Key forms & 12-bar blues (A,A,B); ragtime multi-strain & 32-bar AABA; 12-bar blues; contrafacts & Contrafacts over rhythm changes; 12-bar blues at extreme speed \\
\midrule
Composition vs. improv & Moving from composed (ragtime) toward improvised & Elaborate arrangements (Ellington) vs. head arrangements (Basie) & Minimal arrangement; pre-composed head then extended improv \\
\midrule
Cultural tension & Black innovation, white profit (ODJB) & ``King of Swing'' = white man (Goodman); Cotton Club segregation & WWII anger; Langston Hughes: ``bee-bop'' = police violence \\
\midrule
Key concept & Heterophony, Congo Square, blues codification & Ellington effect, riffs, hot vs. cool, contrafact & Playing the changes, upper extensions, stride piano, walking bass \\
\bottomrule
\end{longtable}
\end{small}

\section{What Each Era Uniquely Contributes}

\subsection{\textcolor{week1color}{Week 1: Early Jazz}}
\begin{enumerate}[leftmargin=1.5em]
  \item \textbf{Invented the genre} --- fused African rhythm, blues, ragtime, European harmony
  \item \textbf{Collective improvisation} --- nobody plays the same line (heterophony)
  \item \textbf{Armstrong's revolution} --- shifted jazz from ensemble to soloist's art; invented scat singing
  \item \textbf{Codified the blues} --- 12-bar A,A,B structure became foundational
  \item \textbf{First recordings} --- spread jazz beyond New Orleans
\end{enumerate}

\subsection{\textcolor{week2color}{Week 2: Swing}}
\begin{enumerate}[leftmargin=1.5em]
  \item \textbf{Big band format} --- Fletcher Henderson's 15-piece template still used today
  \item \textbf{Jazz = America's popular music} --- people danced to it (Lindy Hop)
  \item \textbf{Saxophone rises} --- overtook clarinet as dominant jazz wind instrument
  \item \textbf{Compositional sophistication} --- Ellington elevated jazz to art-music composition
  \item \textbf{Star soloists from ensembles} --- Hawkins, Young, Webster, Hodges
  \item \textbf{Contrafact tradition} --- writing new melodies over existing chord progressions
\end{enumerate}

\subsection{\textcolor{week3color}{Week 3: Bebop}}
\begin{enumerate}[leftmargin=1.5em]
  \item \textbf{Small group revolution} --- WWII economics killed big bands; combos took over
  \item \textbf{Virtuosic improvisation} --- extreme speed, complex harmonic navigation
  \item \textbf{Art music, not dance music} --- deliberate rejection of entertainment model
  \item \textbf{Harmonic expansion} --- upper extensions (9ths, 11ths, 13ths)
  \item \textbf{All 12 keys} --- Parker made this standard for jazz musicians
  \item \textbf{Counterculture identity} --- zoot suits, berets, goatees, horn-rimmed glasses
\end{enumerate}

\section{Key Artists by Era}

\subsection{\textcolor{week1color}{Week 1: Early Jazz}}

\renewcommand{\arraystretch}{1.3}
\begin{small}
\begin{tabularx}{\textwidth}{>{\bfseries}p{3cm} p{3.2cm} X}
\toprule
\textbf{Artist} & \textbf{Nickname/Title} & \textbf{Remember By} \\
\midrule
Scott Joplin & ``King of Ragtime'' & Solo piano, NO improv, syncopated ``oom-pah'' --- the \emph{precursor}, not jazz itself \\
W.C. Handy & ``Father of the Blues'' & First to \emph{write down} blues; tango/``Spanish tinge'' in St. Louis Blues \\
ODJB & --- & First jazz \emph{recording} (1917); all-white band; Keppard refused first \\
Bessie Smith & ``Empress of the Blues'' & Voice bending/sliding into notes; instrumentalists copied \emph{her} \\
Jelly Roll Morton & Creole; ``invented jazz'' & Spoken dialogue intro; funeral march; Red Hot Peppers; composer/arranger \\
Louis Armstrong & ``First Great Soloist'' & Unaccompanied trumpet cadenza; vibrato at END of notes; scat singing; Hot Five/Seven \\
\bottomrule
\end{tabularx}
\end{small}

\subsection{\textcolor{week2color}{Week 2: Swing}}

\begin{small}
\begin{tabularx}{\textwidth}{>{\bfseries}p{3cm} p{3.2cm} X}
\toprule
\textbf{Artist} & \textbf{Nickname/Title} & \textbf{Remember By} \\
\midrule
Duke Ellington & --- & Composer-bandleader; ``Ellington effect''; jungle style; Cotton Club; 50-year band \\
Billy Strayhorn & Ellington's collaborator & Wrote ``Take the A Train'' (NOT Ellington); coined ``Ellington effect'' \\
Count Basie & --- & Minimalist piano; riffs; head arrangements; Kansas City; All-American Rhythm Section \\
Fletcher Henderson & --- & Invented the 15-piece big band format ($\sim$1924) \\
Coleman Hawkins & ``Hawk'' / ``Bean'' / ``Father of Jazz Tenor Sax'' & HOT, vertical improv (arpeggios up/down chord tones); heavy vibrato \\
Lester Young & ``Prez'' (from Holiday) & COOL, horizontal improv (melodic lines across chords); minimal vibrato \\
Ben Webster & (Big Three tenor) & Hot/breathy sound; Ellington's band \\
Benny Goodman & ``King of Swing'' & White clarinetist; integrated his band; Gene Krupa on drums \\
Billie Holiday & ``Lady Day'' (from Young) & Voice like an instrument/horn; fragile timbre; legato phrasing \\
Ella Fitzgerald & ``First Lady of Song'' & Four-octave range; master of scat singing \\
Papa Joe Jones & (Basie's drummer) & Moved four-on-the-floor from bass drum to hi-hat --- influenced bebop \\
\bottomrule
\end{tabularx}
\end{small}

\subsection{\textcolor{week3color}{Week 3: Bebop}}

\begin{small}
\begin{tabularx}{\textwidth}{>{\bfseries}p{3cm} p{3.2cm} X}
\toprule
\textbf{Artist} & \textbf{Nickname/Title} & \textbf{Remember By} \\
\midrule
Charlie Parker & ``Bird'' / ``everyone's God'' & Heart and soul of bebop; alto sax; upper extensions; all 12 keys; heroin addiction \\
Dizzy Gillespie & ``Dizzy'' & Public face of bebop; bullfrog cheeks; upturned trumpet; very high register; jovial \\
Bud Powell & --- & Bebop pianist; right hand = horn-like lines, left hand = sparse comping (not stride) \\
Roy Eldridge & --- & Bridge between Armstrong and Dizzy; hot trumpet sound \\
Charlie Christian & --- & First electric guitar in jazz; died at 25; long flowing melodic lines \\
Art Tatum & --- & Legally blind; incredible technique; stride piano; influenced Bud Powell \\
Jimmy Blanton & --- & Walking bass; died at 23; Ellington's band; huge resonant sound \\
Kenny Clarke & --- & First bebop drummer; drums drive the beat like guitar \\
Max Roach & --- & Featured bebop drummer on Parker's ``Ko Ko'' \\
\bottomrule
\end{tabularx}
\end{small}

\section{Listening Song-to-Era Cheat Sheet}

Use this to instantly place any song in its era. The \textbf{Instant ID Cue} tells you what to listen for.

\subsection{\textcolor{week1color}{Week 1 --- Early Jazz (1900--1928)}}
\emph{If it sounds old, scratchy, small group, collective improv, or stiff/ragtime-ish, it's Week 1.}

\begin{small}
\begin{tabularx}{\textwidth}{p{3.3cm} p{2.8cm} p{0.7cm} X}
\toprule
\textbf{Song} & \textbf{Artist} & \textbf{Year} & \textbf{Instant ID Cue} \\
\midrule
``Maple Leaf Rag'' & Scott Joplin & 1916 & \textbf{Solo piano, mechanical, no swing, no improv} --- ragtime precursor \\
``St. Louis Blues'' & W.C. Handy & 1922 & \textbf{Tango rhythm in the middle} (Spanish tinge); composed; muted trumpet \\
``Dixie Jass Band One-Step'' & ODJB & 1917 & \textbf{Marching band feel}, stiff/military, no real soloists --- earliest jazz recording \\
``Backwater Blues'' & Bessie Smith & 1927 & \textbf{Just voice + piano}; vocal bending/sliding; dual call-and-response \\
``Dead Man Blues'' & Jelly Roll Morton & 1926 & \textbf{Spoken dialogue intro}, then funeral march; Red Hot Peppers \\
``West End Blues'' & Louis Armstrong & 1928 & \textbf{Unaccompanied trumpet cadenza opening} --- most famous moment in early jazz \\
\bottomrule
\end{tabularx}
\end{small}

\subsection{\textcolor{week2color}{Week 2 --- Swing / Big Bands (1925--1945)}}
\emph{If it sounds like a large orchestra with sections (brass vs. reeds), danceable, polished, it's Week 2.}

\begin{small}
\begin{tabularx}{\textwidth}{p{3.3cm} p{2.8cm} p{0.7cm} X}
\toprule
\textbf{Song} & \textbf{Artist} & \textbf{Year} & \textbf{Instant ID Cue} \\
\midrule
``Take the A Train'' & Ellington (Strayhorn) & 1941 & \textbf{Catchy trumpet melody}; Ellington effect underneath; full big band \\
``Cottontail'' & Ellington & 1940 & \textbf{Ben Webster's hot tenor sax solo}; contrafact on ``I Got Rhythm'' \\
``One O'Clock Jump'' & Count Basie & 1937 & \textbf{Melody at the end} (head arrangement); sparse piano; hi-hat timekeeping \\
``Oh! Lady Be Good'' & Basie / Young & 1936 & \textbf{Lester Young's cool, airy, relaxed tenor sax} --- opposite of Hawkins \\
``Body and Soul'' & Coleman Hawkins & 1939 & \textbf{Almost entirely improvised}; melody barely stated; vertical; heavy vibrato; lush \\
``Sing, Sing, Sing'' & Benny Goodman & 1937 & \textbf{Gene Krupa's driving tom-tom drum pattern}; 7+ minutes; clarinet solo \\
``God Bless the Child'' & Billie Holiday & 1942 & \textbf{Fragile, horn-like voice}; legato phrasing; words spill together; yearning \\
``Blue Skies'' & Ella Fitzgerald & 1958 & \textbf{Scat singing showcase}; four-octave range; virtuosic vocal improv \\
\bottomrule
\end{tabularx}
\end{small}

\subsection{\textcolor{week3color}{Week 3 --- Bebop (1940--1950)}}
\emph{If it's blindingly fast, small combo, angular melodies, ride cymbal, and you can't dance to it, it's Week 3.}

\begin{small}
\begin{tabularx}{\textwidth}{p{3.3cm} p{2.8cm} p{0.7cm} X}
\toprule
\textbf{Song} & \textbf{Artist} & \textbf{Year} & \textbf{Instant ID Cue} \\
\midrule
``Dizzy Atmosphere'' & Parker \& Gillespie & Live & \textbf{Extremely fast}; alto sax + trumpet; rhythm changes; Dizzy's high register \\
``Ko Ko'' & Charlie Parker & 1945 & \textbf{Pre-composed intro/outro}; virtuosic solo at breakneck speed; Max Roach; Miles on piano \\
``Parker's Mood'' & Charlie Parker & 1948 & \textbf{The exception}: slow tempo, bluesy, emotional --- proves Parker isn't just about speed \\
``Tempus Fugue-it'' & Bud Powell & 1940s & \textbf{Piano trio}; right hand plays fast horn-like lines; left hand sparse comping (NOT stride) \\
\bottomrule
\end{tabularx}
\end{small}

\section{Quick-Fire Differentiation Tricks}

\subsection{How to Tell the Eras Apart by SOUND Alone}
\begin{enumerate}[leftmargin=1.5em]
  \item \textbf{Can you dance to it?} Yes = Swing. No and it's fast = Bebop. Somewhere in between = Early Jazz.
  \item \textbf{How big is the band?} Full orchestra with sections = Swing. Small combo = Early Jazz or Bebop. Distinguish by speed/complexity.
  \item \textbf{Is everyone improvising at once?} Yes = Early Jazz (heterophony). One soloist over arranged backgrounds = Swing. One soloist blazing over rhythm section = Bebop.
  \item \textbf{How fast and complex?} Moderate and bluesy = Early Jazz. Smooth and danceable = Swing. Blazing fast with angular leaps = Bebop.
  \item \textbf{What does the drummer do?} Marching-band-ish = Early Jazz. Four-on-the-floor bass drum or hi-hat = Swing. Ride cymbal with ``bombs'' = Bebop.
\end{enumerate}

\subsection{How to Tell the Eras Apart by PHILOSOPHY}
\begin{itemize}[leftmargin=1.5em]
  \item \textbf{\textcolor{week1color}{Early Jazz}}: ``Let's all play together and have a good time''
  \item \textbf{\textcolor{week2color}{Swing}}: ``Let's entertain the whole country and get people dancing''
  \item \textbf{\textcolor{week3color}{Bebop}}: ``This is serious art --- if you can't keep up, that's your problem''
\end{itemize}

\subsection{The Throughline of Improvisation}
\begin{itemize}[leftmargin=1.5em]
  \item \textbf{\textcolor{week1color}{Early Jazz}}: Collective improv (everyone at once) $\rightarrow$ Armstrong begins solo improv
  \item \textbf{\textcolor{week2color}{Swing}}: Extended solos within arranged frameworks; hot (Hawkins) vs. cool (Young)
  \item \textbf{\textcolor{week3color}{Bebop}}: Virtuosic solo improv over complex chord changes is the entire point
\end{itemize}

\subsection{The Throughline of the Saxophone}
\begin{itemize}[leftmargin=1.5em]
  \item \textbf{\textcolor{week1color}{Early Jazz}}: Clarinet is the main wind (sax barely present)
  \item \textbf{\textcolor{week2color}{Swing}}: Saxophone overtakes clarinet; Big Three tenors emerge (Hawkins, Young, Webster)
  \item \textbf{\textcolor{week3color}{Bebop}}: Alto sax (Parker) is king; the instrument that defines the era
\end{itemize}

\end{document}
